\section{Harmonic sampling and stopped-interface geometry}
\label{sec:geometry}

We use lattice units in this section. A strand observation specifies the
orientation of finitely many pieces of zero interface, hence only the signs
of adjacent heights. Its conditional law is uniform on the corresponding
convex height polytope. All assertions below concern admissible observations,
that is, observations of positive probability. Constants are uniform in
the number of lattice edges in an observed strand. The number of distinct
macroscopic strands in a comparison is bounded. The exterior heights have
absolute value at most \(1/2\).

The geometric constructions and the random-walk estimates used below are
those of \cite[Sections 2.3 and 3.3--3.7]{SSdiscrete}. We specify their
hypotheses and prove the replacements that concern the law of the field.
In particular, a result for Gaussian interfaces is not used as a theorem
about the present field.

\subsection{Local insertion, signs, and crossings}

\begin{lemma}[Bounded local insertion]\label{g:insertion}
For every fixed \(L<\infty\) there is \(p_L>0\) with the following property.
Condition on any admissible exact oriented strands. In a patch of diameter
at most \(L\), at distance at most \(L\) from one of their tips or from
bounded Dirichlet boundary, prescribe a compatible sign pattern on at most
\(L\) additional vertices. The probability of that pattern is at least
\(p_L\). In particular every feasible prescribed extension of at most
\(L\) interface steps has a probability bounded below by a constant
depending only on \(L\). The assertion remains valid when other positive
string sites have additional upper bounds in \((0,d]\), provided no
prescribed new positive site is one of those upper-bounded sites.
\end{lemma}

\begin{proof}
String heights have absolute value at most one. Lipschitzness therefore
bounds heights in the patch by a deterministic \(M=M(L)\). Enclose it in
a larger fixed box and use Corollary~\ref{f:conditioned-sparse} with a long-distance
deficit greater than \(2M+20\); the small linear-slack version gives this
by enlarging the fixed box. The good-shell event has probability at least
\(1/2\), uniformly in the original strand conditioning.

Fix such a shell and give the new vertices their requested values
\(\pm b\), where \(0<b<1/4\). Clip any feasible filling between the lower
and upper distance cones of these prescriptions. At each other
sign-constrained vertex the clipping interval contains zero, since that
vertex has graph distance at least one from the prescribed vertices.
Consequently clipping preserves its sign. Since the clipping interval
contains zero, it also preserves every old interval \([0,d]\), without
making \(b\) depend on \(d\), and preserves
fixed exterior values of absolute value at most \(1/2\). The shell remains
unchanged because of the long-distance deficit. This is the local
feasible-filling construction underlying Lemma~\ref{f:strict-filling}.

Let \(g\) be the resulting filling. For every original filling \(h\), the
map \(h\mapsto t h+(1-t)g\), with
\(0<t<b/(2(M+b))\), gives the required signs on the patch.
Convexity preserves all the old constraints. If the number of free box
coordinates is \(N\), its Jacobian is \(t^N\); both \(N\) and \(t\) depend
only on the fixed patch and box. Integrating the conditional volume bound
over good shells proves the result. For an interface extension prescribe
the signs on the two sides of its finitely many edges.
\end{proof}

\begin{lemma}[Positive trace on a pure sign side]\label{g:sign-trace}
In a local sector with only positive string boundary, the limiting bounded
harmonic conditional mean has Brownian boundary trace at least \(c>0\),
away from endpoints and incompatible sides. On a negative sector its
trace is at most \(-c\). On an exterior boundary arc alone the trace is
its prescribed value. These statements are sequentially uniform at
diverging lattice scales.
\end{lemma}

\begin{proof}
The conditional law is log-concave and monotone in its boundary data and
in its one-sided constraints. Lemma~\ref{g:insertion} gives a uniform
positive chance that a tested positive string height exceeds some
fixed \(b>0\), even after other positive sites have been restricted to
\((0,d]\). A one-dimensional log-concave density on \([0,1]\) with this
tail bound assigns at most \(C d\) mass to \([0,d]\). One proof integrates
the chord lower bounds for its logarithm between a point in \([0,d]\)
and a point carrying the prescribed tail mass. Thus, for small fixed
\(d>0\), the tested height exceeds \(d\) with conditional probability at
least \(1/2\). Monotonicity makes the conditioning that all earlier
indicators vanish the worst one. Successive uniform couplings show that
the positive string heights dominate \(d\) times independent fair
Bernoulli variables.

Write $A$ for the full original string-height vector. Couple its positive
coordinates with independent fair pins $B_x\in\{0,d\}$ so that
$B_x\le A_x$ on the positive portion, and extend this coupling to the
remaining coordinates of $A$ by disintegration. Conditional on $(A,B)$,
sample the original Gibbs interior given $A$ and a comparison Gibbs
interior pinned to $B$ on the positive portion and to $d/2$ on the
remaining comparison boundary. Then replace this conditional product
by its stationary synchronous joining. Lemma~\ref{c:freezing} preserves
the original field marginal and the comparison mixture. The latter has
the required estimates uniformly in its pins, which lie in
$[0,d]\subset[-1/2,1/2]$. Its law is invariant under
$h\mapsto d-h$, because this map sends the kernel with pins $B$ to
the kernel with pins $d-B$, and the fair product law is unchanged.
Its mean is therefore $d/2$. The filled lower tail of the comparison
has zero fixed input on the positive portion. Theorem~\ref{c:comparison}
and Proposition~\ref{t:varying-trace} give zero trace for every limit
of its mean there; no harmonic equation for the tail is used.
Jensen's inequality bounds the negative part of the limiting mean
difference by this nonnegative zero-trace limit. The two field means are
bounded, harmonic away from the observed sets, and their remote
contribution is bounded by a constant times the harmonic probability of
reaching an endpoint or an incompatible side. This probability tends to
zero along Brownian approaches to the pure portion, by Beurling.
The lower trace is therefore at least \(d/2\). All bounds used are
uniform for each deterministic admissible conditioning sequence, which
gives the claimed sequential uniformity. Reflect the field for the
negative statement.
\end{proof}

\begin{lemma}[The complete-crosscut upper bound]\label{g:upper-half}
For a complete single crosscut, the positive-side conditional mean is
at most \(1/2\), and the negative-side conditional mean is at least
\(-1/2\). This assertion is not made at an open interface tip.
\end{lemma}

\begin{proof}
Condition on all heights on the negative side. At a positive string
vertex, a crossing edge gives an upper bound at most one. Delete the
crossing edges and replace these bounds by independent walls \([0,1]\).
This raises the conditional field by monotonicity. Raising the exterior
positive data to \(1/2\) raises it again. The resulting positive-component
law is invariant under \(h\mapsto1-h\), so its mean equals \(1/2\).
Average the comparison over the conditioned negative heights. The
negative statement follows by reflection. Disconnection of the two
components by the full crosscut justifies deletion of precisely these
crossing constraints.
\end{proof}

\begin{lemma}[High crossings]\label{g:high-crossing}
In a compact corridor of fixed shape and size \(n\), with an observed
component of comparable size and distance, the probability of a path
of heights at least \(K\) crossing a fixed positive fraction of that
corridor tends to zero as \(K\to\infty\), after \(n\to\infty\).
The same is true for the difference of two independent samples under
possibly different bounded exterior or strand conditionings.
In both statements the corridor has an enlarged neighborhood of
width \(cn\), for a fixed \(c>0\), containing no original fixed pins
or sign walls. In the two-field statement this neighborhood lies
in both domains with the same unconditioned lattice graph, and
each original law has the stated macroscopic anchor. The two
conditionings may differ outside this neighborhood.
\end{lemma}

\begin{proof}
We give the two-field proof. Explore every cluster of
\(h_1-h_2\geq K\) connected to a deterministic transverse seed.
Every queried vertex away from the seed has difference at least \(K-2\):
positive queries satisfy the threshold, and failed queries are adjacent
to a positive query. Freeze both samples on the common queried set,
then use the stationary synchronous joining of their conditional laws.
The exploration decisions used only recorded values; hence this
replacement preserves each annealed marginal. All statements below are
annealed over this transcript, not asserted for every possible
conditional pin value.
Use a middle testing subcorridor a fixed positive distance from
the transverse seed. The free-buffer hypothesis removes all
original pins and sign walls from its larger comparison window;
the newly queried common pins there satisfy the displayed
threshold inequality. Thus the one-sided comparison has no
uncontrolled original boundary in that window.

Away from the seed, set
\[
 u_K=\left(1-\frac{h_1-h_2}{K-2}\right)_+ .
\]
It has zero data on the revealed set. The local truncated comparison
estimate in Theorem~\ref{c:comparison}, applied after division by
\(K-2\), gives compact filled extensions with uniformly bounded local
Sobolev energy. Here the cutoff energy grows at most quadratically with
the truncation level: after division it is bounded by the squared cutoff
norm and the macroscopic second moments of \(h_1/(K-2)\) and
\(h_2/(K-2)\). Proposition~\ref{f:averages} controls these moments uniformly.
The same bounds hold after conditioning on an event of probability
bounded below, with its reciprocal as a constant.

Suppose a crossing had probability bounded below along \(K\to\infty\)
and \(n\to\infty\). On that event the revealed set contains a connected
continuum of fixed positive diameter away from the seed. Pass to a
subsequence jointly with this continuum and the filled extensions.
Macroscopic averaged field tightness gives
\((h_1-h_2)/(K-2)\to0\) as a distribution in this iterated limit.
Since \(u_K\geq1-(h_1-h_2)/(K-2)\), its limit is at least one almost
everywhere. The varying-set anchor argument of
Proposition~\ref{t:varying-trace} gives zero trace for that same limit.
Here one uses its samplewise inequalities, without a Poisson equation
for the truncated tail. Choose a probability measure on the crossing
continuum by coordinate projection, with ball masses bounded by $Cr$
in physical units. Move it along the revealed set past each hole to a
retained zero vertex. The maximum hole diameter is polylogarithmic
outside the negligible exceptional event, and the crossing component
cannot be contained in such a hole. Lemma~\ref{t:frostman} bounds the
integral of a radius-$r$ average of the filled tail by
$C(r^\beta+o(1))\|U_n\|_{W^{1,p}}$, for some $\beta>0$.
These inequalities hold for the random measure before expectation.
Joint compactness and strong local $L^p$ convergence permit the mesh
limit at fixed $r$; the bounded conditional Sobolev norm permits
$r\downarrow0$. This proves the zero trace. A \(W^{1,p}\) function which is
at least one almost everywhere cannot have zero trace on a continuum
carrying a fixed positive-diameter Frostman test with exponent
greater than \(2-p\). This is a contradiction.

For one field retain the entire original field on the deterministic
testing window. Given its transcript, compare its unvisited Gibbs
kernel with a symmetric field zero-pinned on the revealed set and
original exterior, and extend this comparison field by zero on its
fixed coordinates. Keep the original field, not its restriction to
a random unvisited component, in the full-window difference.
For each recorded set, deterministic ambient tests of the comparison
field have the bounds of Proposition~\ref{f:averages}; the revealed
crossing component provides a macroscopic anchor on the crossing
event. The original full field keeps its annealed bounds by
disintegration. Queried nonseed differences are at least \(K-1\),
and the same lower tail with \(K-1\) has zero input there. The
free buffer excludes all original walls from this comparison window.
The preceding argument therefore applies without a random mask
on the original field.
\end{proof}

\subsection{Barriers and exact strand comparison}

For a union \(K\) of observed strands, let \(D(K)\) be the domain obtained
by deleting closed lattice triangles meeting \(K\). Write
\[
 d_U(A,B)=\inf\{\operatorname{diam}\alpha:
       \alpha\subset\overline U\text{ joins }A\text{ to }B\}.
\]
The distinction between this internal diameter distance and Euclidean
distance is essential below.

A \((\theta,R)\)-barrier is a path \(\Upsilon\) of diameter in
\((\theta R,R]\), with an observed or exterior boundary component of
diameter at least \(\theta R\) within distance \(\theta^{-1}R\).
Its interior is in \(D(K)\). At each endpoint \(z\) on the boundary,
the first encountered circular crosscut at radius \(\theta R\) bounds
a pocket \(A_z\) with one simple boundary arc; the path meets every
smaller centered circle once inside this pocket. If a point of the
barrier is within \(\theta R/5\) of the observed boundary, it belongs
to one of these pockets and lies within \(\theta R/2\) of its endpoint.
These are the four geometric conditions of
\cite[Section 3.4]{SSdiscrete}. A signed barrier is compatible if the
boundary arc of each endpoint pocket has the indicated sign.
Oppositely signed barriers must have internal distance at least
\(2\theta R\). At bounded lattice scales no hexagon may meet both
prescribed signs.

\begin{proposition}[Uniform barriers]\label{g:barriers}
For every \(\theta>0\) and \(m<\infty\) there is
\(p=p(\theta,m)>0\) such that the following holds at every lattice scale.
Condition on any admissible collection \(K\) of exact oriented strands.
For a compatible collection of at most \(m\) \((\theta,R)\)-barriers,
the probability that all continuations of the observed strands avoid
these barriers is at least \(p\). Endpoints on pure exterior arcs are
allowed provided, for a fixed \(\lambda_*>0\), their positive attachment
arcs have prescribed heights at least \(\lambda_*\), and their negative
attachment arcs have prescribed heights at most \(-\lambda_*\).
For these barriers \(p\) may also depend on \(\lambda_*\).
Interior attachment constants do not depend on \(\lambda_*\).
For diverging-scale applications assume local thickness of the observed
set in the endpoint testing windows: after restriction to a generous
square of side comparable to \(\theta R\) around such a window,
every retained observed component either meets the square's rim or
has diameter at least \(c(\theta)R\). The nearby attachment component
has diameter at least \(c(\theta)R\). Arbitrarily many isolated
microscopic observed islands in these windows are not included.
\end{proposition}

\begin{proof}
We first treat \(R\to\infty\). Away from endpoint pockets the barrier
has clearance \(\theta R/5\). A grid of mesh \(\eta R/100\), with
\(\eta\ll\theta\), meets its bounded thickening in at most
\(C m\eta^{-2}\) cells. Circuits and crossings in those cells form a
fence around the barrier; fill bounded holes in the thickening before
taking its outer fence. Opposite signs have disjoint enlarged supports.
The grid crossings and circuits can be chosen in overlapping rectangles
so that, together with their endpoint attachments, they form a connected
fence. Their number depends on \(\eta,m\), not on the length or
microscopic oscillation of the barrier. We first construct the attachments
to the pure boundary. We will then force the interior fence using
Lemma~\ref{g:high-crossing} and Proposition~\ref{g:shifts}.

Each interior crossing rectangle has a fixed enlargement contained
in the free part of the corridor, by the clearance \(\theta R/5\).
In the endpoint construction only rectangles in the free core are used,
with enlargements smaller than their distance from the removed boundary
neighborhood. Attachment across that neighborhood is the separate
collar problem treated next.

We verify the endpoint attachments, including irregular pockets.
In the annulus with radii \(5\theta R/8\) and \(7\theta R/8\) about
an endpoint \(z\), delete the union \(Q\) of coarse cells meeting the
\(2\eta R\)-neighborhood of the conditioned boundary. Keep the component
containing the barrier as a core. On either side there is a collar
between the core and the pure boundary. Its central testing arc has
diameter at least \(\theta R/32\), distance \(O(\eta R)\) to that
boundary, and access distance at least \(\theta R/32\) to other
conditioned boundary or to either lateral cutoff. These assertions
follow from the radial-pocket and clearance conditions: a path violating
the last assertion would either leave the annulus or cross the barrier
where its clearance is at least \(\theta R/5\). Corner contacts are
resolved by a deterministic perturbation of at most one coarse cell.

Radially select points on the testing arc, move them to centers of
adjacent core cells, and smear them in balls of radius \(\eta R/400\)
inside those cells. For the resulting probability test \(\mu\),
\[
 \mu(B(x,tR))\leq C(t+\eta)/\theta.
\]
For fixed \(\eta,\theta\), all these are smooth tests with uniformly
bounded rescaled norms. The number of coarse cells is bounded, so
subsequential uniformity applies to arbitrary varying old strands.
Beurling and the access separation show that the original mean of
the test is at least \(c+o(1)\) by Lemma~\ref{g:sign-trace}; for
an exterior attachment replace \(c\) by \(c\wedge\lambda_*\).
Its centered variance tends to zero as \(\eta/\theta\to0\) by
Corollary~\ref{t:local-boundary-variance}. Indeed the smearing improves
the displayed ball-mass bound to \(\mu(B(x,tR))\leq Ct/\theta\)
at every scale: for \(t<\eta\), the density bound is
\(C/(\theta\eta R^2)\), giving \(Ct^2/(\theta\eta)\leq Ct/\theta\).
Thus the tests are uniformly Frostman and, for fixed \(\eta\), form
a compact smooth family. The stated local-cut condition supplies the
corollary's thickness hypothesis; no bound on the ambient diameter
divided by \(R\) is imposed.

Explore, within the collar only, the positive cluster attached to its
pure positive boundary. Query a vertex only from an already discovered
positive neighbor. A failed query is negative, and its height is
bounded in absolute value by one. On the negative-crossing event,
the core-facing failed queries contain a negative frontier joining
the two lateral cutoffs and screening the testing cells. The queried
positive vertices lie behind that frontier. This is an adapted
exploration, not a crossing selected after the full field is seen.

Here is the comparison on this event, denoted by $E$. Include all
original string heights in the exact-value transcript; they may be
recorded before the adapted queries. Given the complete transcript,
restrict to the core-facing component cut off by its negative frontier
$F$. The original interior has its exact Gibbs kernel with the
recorded boundary values. As comparison kernel take the uniform law
on this same component with all fixed boundary values replaced by
zero, and denote its field by $h^0$. It is symmetric under height
reflection, so its conditional mean is zero. The zero-pinned laws
have sparse contacts and average bounds uniformly over these component
shapes, using $F$ as a connected anchor. The original estimates remain
annealed estimates under its original strand-conditioned marginal.
Disintegrating and rejoining the two kernels as in
Lemma~\ref{c:freezing} preserves these marginals. It does not require
microscopic slack uniformly for every recorded nonzero pin vector.

For the analytic estimate, retain $h$ on the entire deterministic
testing window, including its coordinates outside the selected
component $C$. Extend only $h^0$ by zero at every vertex outside $C$,
so it is a full ambient zero-pinned field for each recorded shape.
Rejoining may be performed with the full original conditional Gibbs
kernel: its other components retain their conditional baths.
Consequently the full $h$ still has its original annealed marginal.
Put $z=h-h^0$ on this whole window; no component indicator multiplies
$h$. Deterministic ambient zero-mass tests of $h^0$ satisfy
Proposition~\ref{f:averages} uniformly for every recorded shape,
and the coarse coefficient is anchored by $F$. The full original
$h$ has its original anchored and zero-mass bounds. The localized
dyadic argument in \eqref{c:negative} therefore gives
\[
 \E[\|\chi_0 z\|_{H^{-\alpha}}^2\mid E]
 \le C_{\alpha,p_*},\qquad \alpha>0,
\]
when $\Prob(E)\ge p_*$, with the deterministic cutoff $\chi_0$
supported in the fixed larger testing window. Only the original
marginal costs the reciprocal event probability; the zero-pinned
bound is uniform conditional on the recorded shape.

On $F$ the original values are nonpositive. The upper comparison tail
$u=(h-h^0)_+$ consequently has zero fixed input there. Extend it by
zero behind $F$ and apply the one-sided part of
Theorem~\ref{c:comparison}. Use the core-facing components of windows
of a fixed positive fraction of $\theta R$ around the central testing
arc. The lateral cutoffs and any uncontrolled seed cannot be reached
within the corresponding larger window component, by the access
separation above. Thus the cutoff is localized in that component;
all edges to the removed positive cluster pass through $F$.
The nearby original boundary anchors the original marginal, and
$F$ anchors the comparison marginal. The zero-extended tail is
uniformly Lipschitz across the frontier. If its hole extension has
a positive copied value $U(x)=u(y)$, the retained rim site $y$ lies
in $C$ and $u(y)=z(y)>0$. Full-window Lipschitzness of $z$ gives
\[
 U(x)^2\le z(x)U(x)+C(1+d_x)U(x).
\]
At $U(x)=0$ this is automatic. The inequality remains valid even
if a hole copies a positive value into the removed side; it does
not require an exact-zero filled extension on every internal
component. Pairing it with a deterministic cutoff uses the full
$z$ in Sobolev duality, as in \eqref{c:one-sided-pointwise}--
\eqref{c:one-sided-absorb}, rather than a randomly masked original
field. Retained frontier vertices still have value zero.
Thus the expected energy and
squared $W^{1,p}$ norm of the filled tail $U$, for $3/2<p<2$, have
bounds independent of $\eta$. Restricting to $E$ of probability at
least $p_*>0$ costs only a factor $p_*^{-1}$ in these annealed bounds.

The testing measure is within $C\eta R$ of $F$: a testing cell has
a path of that diameter to its nearest pure boundary piece, and the
screening frontier must intersect this path. Transport its mass to
$F$ and then past the polylogarithmic holes to retained zero vertices.
The frontier is connected between the lateral cutoffs and has diameter
comparable to $\theta R$, so this transport has the common Frostman
bounds. The samplewise averaging and anchor estimates of
Lemmas~\ref{t:frostman}--\ref{t:anchors}, in units of $\theta R$,
give, for some $\beta>0$,
\[
 \limsup_{R\to\infty}
 \E\bigl[|\langle\mu,U\rangle|^2\mid E\bigr]
 \le C_{p_*}(\eta/\theta)^{2\beta}.
\]
All geometry and $\eta>0$ are fixed before this lattice-scale limit.
The original--filled $L^1$ estimate transfers the tested first moment
at that fixed scale. Since $h\le h^0+u$ and the comparison kernel has
zero mean conditional on the transcript,
\[
 \limsup_{R\to\infty}\E[\langle\mu,h\rangle\mid E]
 \le C_{p_*}(\eta/\theta)^\beta.
\]
This is not a nonpositive bound at fixed $\eta$; its right side vanishes
only afterwards as $\eta/\theta\downarrow0$. No Poisson equation or
Poisson uniqueness theorem is applied to the comparison tail.
For a variable \(X\) with variance \(v\),
\[
 |\E[X\mid E]-\E X|\leq\sqrt{v/\Prob(E)}.
\]
Thus the collar-crossing probability tends to zero. The reflected
argument treats a negative pocket.

It remains to join these attachments by the interior fence. Choose
\(\eta\) so that the collar failures together have probability less
than \(1/10\) for large \(R\). In each positive fence rectangle,
planar crossing duality gives a path with heights greater than \(-H\)
unless there is a transverse path with heights at most \(-H\).
Apply Lemma~\ref{g:high-crossing} to the reflected field to exclude
the latter event. The reflected construction supplies paths with
heights less than \(H\) in negative rectangles. The same crossing
construction in the finitely many overlapping rectangles supplies
any required circuits. Choose fixed \(H\) sufficiently large that
all these threshold-crossing failures together have probability less
than \(1/10\). The successful threshold crossings and collar attachments
therefore occur together on an event \(E_{\rm fence}\) of probability
at least \(4/5\).

The prescribed shift is deterministic given the conditioned geometry:
set it equal to \(H+1\) on each positive crossing rectangle and to \(-(H+1)\)
on each negative crossing rectangle. Interpolate to zero within their
fixed free enlargements, keeping zero at all observed sites. The separated supports give a bounded
prescription with edge differences at most \(C(H,\theta,\eta)/R\),
as required by Proposition~\ref{g:shifts}. Keep its sign reinforcing
on the corresponding collar. The corrected shift differs from this
prescription by \(o(1)\), so it turns each threshold path into a
strict sign path. It also preserves the successful collar attachments:
the compatibility and clearance conditions give internal diameter
distance at least \(cR\) from a collar or prescription support of
one sign to a prescription support of the other. The near-contact
components used by the shift construction have diameter \(o(R)\)
in the free graph. A component meeting a positive collar therefore
contains no negative prescription, and its averaged
shift is nonnegative; the negative case is reflected. Grounded
components have zero shift.

All of \(\theta,\eta,H\) are fixed before \(R\to\infty\).
Apply Proposition~\ref{g:shifts} once to \(E_{\rm fence}\). It gives a uniformly
positive probability for the connected sign fence under the original
law. That fence prevents every continuation from entering the protected
side. This proves the lower bound above a fixed threshold
\(R_0(\theta,m)\).

For \(R\leq R_0\), the barrier thickening and the required finite sign
fences lie in bounded patches at bounded distance from the anchoring
components. Lemma~\ref{g:insertion} forces them with probability
bounded below. Taking the smaller of the two bounds proves the
all-scale statement.

Here is the local-thickness check for the applications below. External
arms and adapted continuations are attached to the original exterior.
An internal record is a connected incident strand through the sampling
center, of diameter comparable to its record radius when its quality
is positive. An endpoint testing square can be chosen small enough
not to contain that center. If an internal record of smaller diameter
meets such a square, its connection to the center crosses the square's
rim. Otherwise its diameter is already comparable to the square's
scale. In the extra-excursion construction the endpoint squares lie
near radius \(\rho\), have size at most \(c\rho\), and are disjoint
from the center even when the inner record radius is much smaller
than \(\rho\). Thus no global lower bound on all observed strand
diameters is being presumed there. In the strand-swap application
the comparison annulus is free of all the inner observations.
\end{proof}

\begin{lemma}[Exact-strand comparison]\label{g:strand-comparison}
Fix concentric inner and outer disks with radii in a fixed ratio
strictly greater than one. Consider two base conditioning systems
whose domains contain the closed outer disk and have the same
unconditioned lattice graph there. All vertices constrained by the
original strand observations, additional walls, or Dirichlet data
must be outside the outer disk in both systems. Thus the entire outer
disk, not just its annulus, is free in the two base laws. Each system
has anchored observed or Dirichlet pieces of comparable diameter
and distance outside this disk. Let \(I\) specify a bounded collection
of oriented strands all of whose adjacent sign-constrained vertices
lie in the inner disk. Then
\[
 C^{-1}\Prob_2(I)\leq\Prob_1(I)\leq C\Prob_2(I),
\]
with a constant depending only on the fixed ratios and number of
strands. The graphs need agree only inside the outer disk.
The assertion does not compare laws with different pre-existing
inner observations.
\end{lemma}

\begin{proof}
Take two independent samples from these base laws, conditioning the
first additionally on \(I\). The only observations inside the outer
disk are now this additional event in the first replica, and their
support lies inside the inner disk. Hence the annulus remains free
in both replicas. Each corridor used there has an anchor at comparable
distance under each law; an inner strand from \(I\) is also an
allowed anchor for the first replica.
Choose the two circuit supports in a closed middle subannulus
and cover them by finitely many corridors whose fixed
enlargements remain in that annulus. They satisfy the
free-buffer hypothesis of Lemma~\ref{g:high-crossing} in
both replicas. Lemmas~\ref{g:high-crossing}
and \ref{g:shifts} therefore give a
uniformly positive probability of two nested circuits in the free
annulus on which \(h_1-h_2\) has opposite strict signs.
We use the cluster-swapping construction of Sheffield
\cite[Chapter~8]{SheffieldSurfaces}, in the hard-constraint form
described below; see also \cite[Section~2]{CAP}.
For each edge declare it a dependency edge when swapping the
replicas at just one endpoint violates an edge constraint in one
of the replicas. This graph depends only on the unordered height
pairs. If the order of the replicas reverses across an edge, either
partial swap is feasible: each mixed endpoint difference lies
between differences of the original feasible samples. Hence a
dependency path cannot change the sign of \(h_1-h_2\).
No dependency component meeting the inner disk reaches outside
the two opposite-order circuits.

Swap both replicas at every vertex in a dependency component
meeting the inner disk. Every swapped component stays inside the
outer disk. The base laws impose no vertex observations there;
all their original observations therefore remain at untouched
vertices. A graph edge across a swapped-component boundary is
not a dependency edge, so the partial swap on that edge is feasible.
This proves preservation of both base polytopes. All vertices
tested by \(I\) are swapped, which transfers \(I\) from the first
replica to the second; the first replica is not required to retain
this additional event. The operation is an involution because the dependency
graph and selected components are invariant under swaps.
On each component-selection region it is a coordinate permutation,
so it preserves Lebesgue measure. If \(V_i\) is the volume of the
polytope with exterior conditioning \(i\), the uniform positive
fraction of swappable pairs gives
\[
 c\,V_1(I)V_2\leq V_1V_2(I).
\]
Reverse the experiments for the other inequality.
\end{proof}

\subsection{The harmonic-landing records}

Use rounded disks \(\mathfrak B_R\) made of dual hexagons.
Fix the origin and an incident dual edge \(e_\sigma^*\). Let \(S\)
be independent simple random walk from \(v\), let \(\tau_0\) be
its first visit to zero, and reverse its portion up to \(\tau_0\)
to obtain \(\breve S\). Let \(\tau\) be the first visit to a
vertex adjacent to the completed chord or to exterior boundary.
Define
\[
 Z^\sigma=\{S_\tau=0,\ e_\sigma^*\subset\gamma\}.
\]
Let \(\operatorname{ext}_R\gamma\) consist of the two initial
pieces grown from the marked endpoints until first entering
\(\mathfrak B_R\); it is the whole chord if the disk is missed.
Let \(\beta_R^\sigma\) be the part in \(\mathfrak B_R\) of the
zero interface through \(e_\sigma^*\), connected to that edge,
and put it equal to the empty path if there is no such interface.
Let \(\operatorname{ext}_R S\) end at the first entrance to
\(\mathfrak B_R\), and let \(\operatorname{int}_R\breve S\)
end at its first exit. Set
\[
 \Phi_R=(D,\partial_+,h_\partial,v,
              \operatorname{ext}_R\gamma,\operatorname{ext}_R S),
 \qquad
 \Theta_R=(\beta_R^\sigma,\operatorname{int}_R\breve S).
\]
The geometric path in \(\Theta_R\) is initially unoriented.
Conditioning on its orientation means specifying the adjacent signs.
Every subsequent field-law application below conditions on these
exact signs, then averages over the orientations.

Let \(x,y\) be the two external strand tips and \(q\) the external
walk tip at radius \(R\). When there has been no forbidden contact,
define
\[
 Q(\Phi_R)=\frac1{100}\wedge
 \frac{
 d(x,\operatorname{ext}_R S)\wedge
 d(y,\operatorname{ext}_R S)\wedge
 d(q,\operatorname{ext}_R\gamma)\wedge |x-y|}{R}.
\]
Put \(Q=0\) if the disk is missed by the chord or the revealed walk
hits a revealed strand. Define \(Q(\Theta_R)\) by the same formula
using the two inner strand tips and reversed-walk tip; put it equal
to zero for a closed inner interface or a forbidden walk contact.

Figure~\ref{fig:harmonic-records} depicts the two records and the
unexposed annulus between them.

\begin{figure}[htbp]
\centering
\begin{tikzpicture}[scale=.8,>=Stealth,font=\small]
  \fill[gray!10] (0,0) circle (3);
  \fill[white] (0,0) circle (1.15);
  \draw[dashed,gray] (0,0) circle (3);
  \draw[dashed,gray] (0,0) circle (1.15);
  \draw[very thick,blue!65!black]
    (.15,3.65) .. controls (-.1,3.4) and (.35,3.2) .. (.2,3);
  \draw[very thick,blue!65!black]
    (.1,-3.65) .. controls (.4,-3.4) and (-.15,-3.2) .. (.05,-3);
  \draw[very thick,blue!65!black]
    (.03,-1.15) .. controls (-.1,-.65) and (.2,-.4) .. (.15,.05)
    .. controls (.1,.5) and (.35,.8) .. (.16,1.14);
  \draw[very thick,red!65!black,->]
    (-3.9,.3) .. controls (-3.6,.1) and (-3.3,.3) .. (-3,.1);
  \draw[very thick,red!65!black,->]
    (0,0) .. controls (-.4,-.15) and (-.7,.18) .. (-1.14,.12);
  \fill (0,0) circle (.045);
  \node[below left] at (0,0) {\(0\)};
  \node[above left] at (-3.7,.32) {\(v\)};
  \node[anchor=west] at (1.35,.95) {\(\Theta_r\): inner record};
  \draw[gray] (1.3,.9)--(.65,.6);
  \node[anchor=west] at (2.4,2.65) {\(\Phi_R\): outer record};
  \draw[gray] (2.35,2.65)--(.35,3.2);
  \node at (1.8,-1.65) {buffer};
  \node[anchor=east] at (-.95,-1.8) {sampled side};
  \draw[gray,->] (-1.05,-1.55)--(-.45,-.4);
  \node[above right] at (.77,.77) {\(r\)};
  \node[above right] at (2.1,2.1) {\(R\)};
  \draw[very thick,blue!65!black] (-2,-4.15)--(-1.45,-4.15);
  \node[anchor=west] at (-1.3,-4.15) {observed interface};
  \draw[very thick,red!65!black] (-2,-4.55)--(-1.45,-4.55);
  \node[anchor=west] at (-1.3,-4.55) {observed walk pieces};
\end{tikzpicture}
\caption{The two records leave an annulus available for separation
and hookup. The inner walk arrow indicates reversal from the harmonic
landing. The shaded annulus is unexposed; no interface connection there
has yet been required.}
\label{fig:harmonic-records}
\end{figure}

\begin{lemma}[Hookup probability]\label{g:hookup}
For \(9R/8<R'<5R\), with the observer and exterior boundary outside
\(\mathfrak B_{6R}\), there is a universal upper bound for all
\(R\geq R_0\), where \(R_0>1\) is fixed:
\[
 \Prob(Z^\sigma\mid\Phi_{R'},\Theta_R)\leq C/\log R.
 \tag{G1}
\]
For every \(a>0\) there is \(c(a)>0\) such that the same probability
is at least \(c(a)/\log R\) when
\(Q(\Phi_{R'})\wedge Q(\Theta_R)\geq a\) and \(R\geq R_0(a)\).
\end{lemma}

\begin{proof}
The field part of \(\Phi_{R'}\) constrains only sites adjacent to
the external arms, hence outside radius \(R'-O(1)\).
The sites specifying \(\beta_R^\sigma\) lie inside radius \(R+O(1)\).
Choose comparison disks between these two radii, with
outer radius \((R+R')/2\) and inner radius \(R+(R'-R)/4\).
For sufficiently large \(R\), they have a free intervening annulus
of fixed relative width. The original base observations are all
outside the outer disk, and the external arms give the required
anchors. Thus Lemma~\ref{g:strand-comparison} applies to the two
orientations of the inner strand, comparing the exterior data with
their global sign reversal. Together with sign symmetry it gives a uniformly
positive probability that the two observed unoriented configurations
have compatible signs on the side containing the walk. In that case
join corresponding interface tips and the two walk tips by three
disjoint corridors, of width \(c(a)R\). Four signed barriers alongside
the two interface corridors have the required pure endpoint pockets.
Proposition~\ref{g:barriers} forces the interface joins within these
corridors with uniformly positive probability. Planarity forces the
outer arms to merge with the corresponding inner arms; there is no
other exit through the fences.

The remaining walk bridge must reach the joining vertex of the reversed
inner piece. Its conditioning to reach this vertex before zero has
probability bounded below. Harnack gives a positive probability to
follow the free third corridor to a fixed relative neighborhood of
that vertex; the two-dimensional point-hitting estimate gives
probability comparable to \(1/\log R\) to hit it there. The loop portion
between its first and last visits stays in that corridor with probability
bounded below, by the last-exit decomposition. This proves the lower
bound. Conditional on any completed field configuration, hitting the
same vertex before a nearby interface continuum has probability at
most \(C/\log R\); this proves the upper bound. These are precisely
the walk hypotheses of \cite[Lemma 3.14]{SSdiscrete}: radii in the
specified compact ratio range, observer outside the larger disk,
and two prescribed first/last walk pieces. All field estimates used
in its corridor construction have been supplied above.

We record the two augmented upper bounds used below. First, the upper
bound still holds after the whole field, hence \(\gamma\), is fixed;
the middle walk bridge remains unexposed. Second, put
\(H=(\Phi_{3r},\Theta_{2r})\), and expose that middle walk only until
its first exit from a circle of radius \(\rho/3>10r\), if this occurs
before its prescribed terminal piece. Call this stopped prefix
\(S_{**}\). On this escape event,
\[
 \Prob(Z^\sigma\mid H,S_{**},\hbox{escape})
       \leq C/\log r.                              \tag{GK}
\]
For verification, let \(\hat q\) be the joining vertex of the fixed
reversed inner piece. The first-hit/last-exit decomposition leaves a
walk from the current exit vertex to \(\hat q\), conditioned to hit
\(\hat q\) before zero, followed by its last-exit loop. Harnack bounds
the conditioning probability below by a fixed constant, since the
exit vertex is outside \(B_{10r}\) while
\(|\hat q|=2r+O(1)\). To realize \(Z^\sigma\), this first part
must hit \(\hat q\) before the connected inner interface. The same
point-versus-continuum estimate as above is at most \(C/\log r\);
the remaining loop can only reduce the probability. This argument
holds for every completed field consistent with the records, and
averaging gives (GK). It does not expose the later successful bridge.
\end{proof}

\begin{lemma}[Separation under a harmonic landing]\label{g:separation}
For every \(\eps>0\) there is \(c(\eps)>0\) such that, for all
sufficiently large \(R\),
\[
 \Prob\bigl(
 Q(\Phi_{3R})\wedge Q(\Theta_{2R})>c(\eps)
 \mid \Phi_{4R},\Theta_R,Z^\sigma\bigr)>1-\eps
 \tag{G2}
\]
uniformly in all admissible records for which the conditioning has
positive probability. The observer and original boundary are outside
\(\mathfrak B_{6R}\).
\end{lemma}

\begin{proof}
First suppose the inner quality is bounded below. Expose the external
arms inward, taking \(r_{j+1}=(1-10Q_j)r_j\), where \(Q_j\) is their
quality or zero if the hookup has become impossible. Conditional on
the current records, but not on \(Z^\sigma\), there are constants
\(c_0,c_1>0\) such that
\[
 \Prob(Q_{j+1}\geq(2Q_j)\wedge 1/100\mid\Phi_{r_j},\Theta_R)
       \geq c_0,
 \tag{G3}
\]
\[
 \Prob(Q_{j+2}=0\mid\Phi_{r_j},\Theta_R,Q_{j+1}<2Q_j)
       \geq c_1. \tag{G4}
\]
For (G3), steer the two interface tips and the walk into three separated
corridors of size \(r_jQ_j\). Their annulus is disjoint from the fixed
inner observations. Proposition~\ref{g:barriers} and the local Harnack
bound for the walk bridge give the required probability whenever
\(r_jQ_j\) exceeds a fixed threshold. Below that threshold the finite
step construction of Lemma~\ref{g:insertion} applies. The steering
sites lie within a bounded distance of the existing tips.

For (G4), put \(\mathcal F_j=\sigma(\Phi_{r_j},\Theta_R)\),
\(q=Q_j>0\), \(r'=r_{j+1}\), and
\(F=\{Q_{j+1}<2q\}\). Reveal the two external arms to radius \(r'\),
but do not yet reveal the whole intervening walk. Write \(x,y\) for
their tips, \(w\) for the eventual walk tip at radius \(r'\), and
\(S^\dagger\) for the walk between its first radius-\(r_j\) and
radius-\(r'\) entrances. The following five events cover \(F\):
\[
\begin{split}
 M_0&=\{Q_{j+1}=0\},\\
 M_1&=\{d(S^\dagger,\{x,y\})<(2q r')\wedge|x-y|\},\\
 M_2&=F\cap\{r'Q_{j+1}=d(w,\operatorname{ext}_{r'}\gamma)\},\\
 M_3&=F\cap\{0<r'Q_{j+1}=|x-y|\},\\
 M_4&=F\cap\{Q_{j+1}=1/100\}.
\end{split}                                                    \tag{G4a}
\]
An event involving tips is empty if those tips do not exist; this
case is already in \(M_0\). The old walk outside radius \(r_j\)
cannot cause the new insufficient separation, since
\(r_j-r'=10q r_j>2q r'\), with the rounded-disk lattice margin
absorbed in these constants. Thus \(M_i\subset F\) and their union
is \(F\), including ties in the quality minimum.

Let \(L=\{Q_{j+2}=0\}\). The field-free stopped-walk estimate in
\cite[proof of (3.38), case \(i=1\)]{SSdiscrete} gives
\[
 \Prob(L\cap M_1\mid\mathcal F_j)
       \geq c\Prob(M_1\mid\mathcal F_j).             \tag{G4b}
\]
Here are the conditioning and geometric hypotheses for that estimate.
After fixing \(\operatorname{ext}_{r'}\gamma\), stop \(S^\dagger\)
at its first visit within
\(a=(2q r')\wedge|x-y|\) of either tip. Reveal only the walk up to
this first approach. The conditioned inner reversed record leaves
the first walk bridge as a simple walk transformed to hit its
joining vertex before zero. Harnack compares this transform with
ordinary walk by fixed factors in the approach neighborhood: the
fixed inner record is at distance comparable to \(R\), whereas
\(a\leq2q r'\leq r'/50\). The approached strand extends a distance
comparable to \(a\), the other tip is at distance at least \(a\),
and the preceding annular width is at least \(5a\).
These are exactly the hypotheses of the source's first-approach
walk argument. Its strong Markov and planar hitting estimate have
a numerical lower bound, independent of \(a/r_j\), and force a
strand hit before the next quality-controlled annulus is passed.
Integrating this stopped-record estimate gives (G4b). In particular
one must not additionally fix the later walk segment after it has
escaped from this close approach; a uniform estimate for each such
completed record would be false.

For \(M_2\) and \(M_4\) the same hitting bound starts at the current
walk tip and does apply after that radius-\(r'\) record is fixed.
On \(M_3\), a compatible barrier around the midpoint of \([x,y]\)
at radius \(2|x-y|\) forces the two strands to join before the next
inner circle, since its radial gap is \(10|x-y|\). This kills hookup
without meeting the fixed inner strand. Proposition~\ref{g:barriers}
applies, and bounded scales use Lemma~\ref{g:insertion}. The case
\(M_0\) is immediate. Hence (G4b) holds for every \(M_i\), with a
common \(c>0\). The events need not be disjoint: summing yields
\[
 5\Prob(L\cap F\mid\mathcal F_j)
 \geq\sum_{i=0}^4\Prob(L\cap M_i\mid\mathcal F_j)
 \geq c\Prob(F\mid\mathcal F_j).
\]
Division by \(\Prob(F\mid\mathcal F_j)>0\) proves (G4).

Iterating (G3), followed by Lemma~\ref{g:hookup}, shows that a quality
\(q>0\) with a fixed annular buffer remaining has
\[
 \Prob(Z^\sigma\mid\hbox{current records})
       \geq c q^\alpha/\log R                         \tag{G5}
\]
for a fixed \(\alpha>0\). The radial loss while doubling the quality
is bounded by a geometric sum. Let \(M_n\) count visits to
\((2^{-n},2^{1-n}]\). By (G4), survival for \(2m\) such visits has
unconditioned probability at most \((1-c_1)^m\).
Multiply by the upper bound in (G1), and divide by (G5) at the first
visit to the band. This yields
\[
 \Prob(M_n>2m\mid\Phi_{4R},\Theta_R,Z^\sigma)
       \leq C c_0^{-n}(1-c_1)^m.                     \tag{G6}
\]
For a sufficiently large fixed \(A\), the bounds in (G6) with
\(m=An\) have a geometric tail in \(n\). Choose \(N\) so that their
sum over \(n\ge N\) is less than \(\eps/3\), and so that
\[
 C A\sum_{n\ge N}n2^{-n}<1/2.
\]
Here \(C\) bounds the radial loss: one visit to the \(n\)-th band
costs at most \(10\cdot4R\cdot2^{1-n}\), since
\(r_j-r_{j+1}=10Q_jr_j\). Outside the exceptional event in (G6),
there are at most \(2An\) such visits. Thus, while the quality stays
below \(2^{-N}\), the total radial loss is less than \(R/2\).
With probability at least \(1-\eps/3\), a fixed positive quality is
therefore attained before the radius falls from \(4R\) to \(7R/2\).

To rule out a subsequent collapse at \(3R\), use the two mechanisms
in \cite[proof of Lemma~3.15, equations (3.40)--(3.42)]{SSdiscrete}.
Condition on the first record at which that positive quality is attained,
and reveal the external strands to radius \(3R\). If their tips are within \(sR\), their midpoint is
now fixed. Expose the continuations successively to barriers at radii
\(2sR,4sR,\ldots,cR\) about that midpoint. Hookup requires
breaching each barrier. Proposition~\ref{g:barriers} bounds each
conditional breach probability by \(1-p\); the number of scales is comparable to
\(\log(1/s)\), giving
\(Cs^\beta\) for a fixed \(\beta>0\). Fixing the completed field while leaving the walk bridge
unexposed supplies the additional upper factor \(C/\log R\) from (G1).
For a walk--strand approach, fix \(\operatorname{ext}_{3R}\gamma\)
and reveal the walk only to its first \(sR\)-approach to that strand.
Beurling gives \(Cs^\zeta\), for a fixed \(\zeta>0\), for reaching radius
\(5R/2\) without a hit. At its first entrance there, the remaining
unexposed bridge again contributes \(C/\log R\). Hence the joint probability of late collapse and hookup
is at most \(C(s^\beta+s^\zeta)/\log R\). Divide by (G5) at the
already attained positive quality, then take \(s\) small.

For general inner quality expose both records: increase the inner
radius by factor \(1+10Q_j\) while decreasing the outer radius by
\(1-10Q_j\), using the smaller of their qualities. Until the target
radii \(2R\) and \(3R\), their annuli are disjoint. The same finite
collection of compatible barriers proves (G3)--(G4) for this minimum;
independence of the fields in the two annuli is not needed.
The same band calculation and late-collapse bound establish (G2).
The conditional records in these steps are exactly the two-sided
exploration records used in \cite[Lemma 3.15]{SSdiscrete}; their
field part consists solely of specified oriented strands.
\end{proof}

\begin{proposition}[Comparison of sampled configurations]
\label{g:sample-comparison}
There are \(C<\infty\) and \(R_0<\infty\) such that the following holds
for \(R\geq R_0\). Take two domains
containing \(\mathfrak B_{6R}\), with observers outside that disk,
and any two external records \(\phi,\phi'\) at radius \(4R\) having
positive hookup probability. For every event \(H\) determined by
\(\Theta_R\), with the same specified incident-edge label and landing
side in the two experiments,
\[
 C^{-1}\Prob'(H\mid\Phi'_{4R}=\phi',Z'^\sigma)
 \leq \Prob(H\mid\Phi_{4R}=\phi,Z^\sigma)
 \leq C\Prob'(H\mid\Phi'_{4R}=\phi',Z'^\sigma).
 \tag{G7}
\]
The constant is independent of \(R\) and of additional unused annuli
outside radius \(4R\).
\end{proposition}

\begin{proof}
Before hookup conditioning, the field observations in each
\(\Phi_{3R}\) are only the signs adjacent to its two external
arms, supported outside radius \(3R-O(1)\). The event specifying
an oriented \(\beta_{2R}^\sigma\) tests vertices inside
radius \(2R+O(1)\). Apply Lemma~\ref{g:strand-comparison} with
inner disk of radius \(9R/4\) and outer disk of radius \(5R/2\).
Both base laws have no observations in this outer disk, their graphs
agree there, and their external arms connect scale \(3R\) to the
original boundary outside scale \(6R\), providing the anchors.
Thus the inner oriented-strand laws are comparable by a fixed factor
under these different external records. At this step no
\(\Theta\)-record has already been imposed in either base law:
the inner record is precisely the additional event being compared.
The reversed-walk law is also comparable by a
fixed factor: a walk from either entrance point has probability
bounded below to visit the other before reaching zero. Apply the
strong Markov property at that visit. The field and walk factors
are independent before hookup. Thus there is a reference weight
\(p_{2R}(\theta)\) for inner records such that their probabilities
given a record at radius \(3R\) are comparable to this weight.

Choose the quality threshold in Lemma~\ref{g:separation} with error
\(1/2\), and let \(X\) be the inner records at \(2R\) exceeding it.
Let \(\mathcal Q\) be the event that both the \(2R\) inner and
\(3R\) outer records exceed it. Separation holds even conditional
on a specified \(\Theta_R=\vartheta\). Hence conditioning further
on \(\mathcal Q\) changes the probability of this specified inner
record by at most a fixed factor. On \(\mathcal Q\), Bayes' formula,
the reference-weight comparison, and Lemma~\ref{g:hookup} give
\[
 \Prob(\Theta_{2R}=\theta\mid\mathcal Q,Z^\sigma,\Phi_{3R})
 \asymp \frac{p_{2R}(\theta)}{\sum_{\xi\in X}p_{2R}(\xi)}
 \quad(\theta\in X).
\]
Indeed the numerator before normalization is comparable to
\(p_{2R}(\theta)/\log R\), and summing determines its normalization.
Average over \(\Phi_{3R}\), then sum over all \(\theta\in X\)
whose restriction to radius \(R\) is \(\vartheta\).
The resulting reference ratio is independent of the external
experiment, proving (G7) first for single records, then for events.
Conditioning on data farther out and averaging over the radius
\(4R\) records preserves the same constant.
\end{proof}

\begin{lemma}[No additional sampled excursions]\label{g:extra}
Let \(\mathcal J(r,R)\) be the event that more than two disjoint
interface arcs cross from \(\mathfrak B_r\) to
\(\partial\mathfrak B_R\), or that the walk exits
\(\mathfrak B_R\) after first entering \(\mathfrak B_r\) and before
its landing at zero. For every \(\eps>0\) there are
\(A(\eps)<\infty\) and \(r_0(\eps)<\infty\)
such that
\[
 \Prob(\mathcal J(r,R)\mid\Phi_R,\Theta_r,Z^\sigma)<\eps
 \tag{G8}
\]
when \(r\geq r_0(\eps)\), \(R>A(\eps)r\), the observer and boundary are outside
\(\mathfrak B_{4R}\), and the conditioning has positive probability.
The same holds, with a harmless change of \(\eps\), conditioned
only on \(\Phi_R\) and \(Z=\{S_\tau=0\}\).
\end{lemma}

\begin{proof}
Separation in inner and outer buffers makes all four records at
scales \(2r,3r,R/2,3R/4\) have quality at least \(c(\eps)\), except
with probability \(O(\eps)\). Put \(\rho\asymp\sqrt{rR}\).
Let \(H=(\Phi_{3r},\Theta_{2r})\). The walk portion \(S_*\) from its
first radius-\(3r\) entrance to the last visit, before zero, to the
endpoint of its radius-\(r\) reversed piece is unlikely to leave
\(\mathfrak B_{\rho/3}\) without hitting the external arms. The
ordinary Beurling estimate and first-hit/last-exit decomposition
bound this probability, conditional on \(H\), by
\(b(R/r)\to0\). If it escapes, reveal only its prefix \(S_{**}\)
up to first exit. The augmented bound (GK), not a bound conditioned
on the completed walk, gives
\[
 \Prob(S_*\not\subset\mathfrak B_{\rho/3},Z^\sigma\mid H)
       \leq C b(R/r)/\log r.
\]
Divide by the lower bound \(c(\eps)/\log r\) on separated records
and average over them. This proves confinement under the landing law.

For the outer construction, use only \(\operatorname{ext}_{3r}S\).
Let \(S'\) be its terminal path from its last exit of
\(\mathfrak B_{R/3}\) to its endpoint at radius \(3r\).
Choose measurably from \(S'\) an arc on the circle of radius \(\rho\)
whose union with \(S'\) separates the inner and outer circles.
Intersect this open circular arc with \(D(\operatorname{ext}_{3r}\gamma)\).
Its two components next to \(S'\) have other endpoints on the boundary
of this deleted-triangle domain, not on the interface itself;
call them \(\alpha_1,\alpha_2\). Thus an endpoint attachment does
not already count as an interface intersection. They are
\(\Phi_{3r}\)-measurable. In the component on either side, place
nested barriers separating such a piece from the inner circle.
The deterministic construction is the clearance construction in
\cite[proof of (3.48)]{SSdiscrete}: take a separating circular arc;
if its middle has adequate boundary clearance, attach shortest
endpoint segments, and otherwise repeat at the smaller distance
between its two facing pure boundary pieces. The procedure
terminates because the two boundary pieces are disjoint.
Each output has a fixed relative barrier parameter, although its
lattice diameter may be small. Barriers selected in distinct
dyadic annuli remain mutually separated.

After exposing the two interface arms to one barrier, the next
barrier satisfies Proposition~\ref{g:barriers}, including its
bounded-scale case. Thus crossing \(N\) barriers has probability
at most \((1-c)^N\), conditional on \(H\) but not on the landing.
For \(F=\{\gamma\cap(\alpha_1\cup\alpha_2)\ne\varnothing\}\),
fix the completed field while retaining the unexposed walk bridge.
The first augmented upper bound following (G1) gives
\[
 \Prob(F,Z^\sigma\mid H)
 =\E[1_F\Prob(Z^\sigma\mid H,\gamma)\mid H]
 \leq \frac{C}{\log r}\Prob(F\mid H).
\]
The same lower normalization gives \(C(\eps)(1-c)^N\) under the
landing law, where \(N\to\infty\) with \(R/r\). Repeat the whole
construction from the inner reversed record \(\Theta_{R/3}\),
without also revealing the forward intermediate bridge.

Outside these exceptional events, confinement identifies the terminal
forward path \(S'\) with the corresponding portion of the reversed
inner record. The same deterministic circular-arc selection then
makes the two crosscuts from the separate descriptions coincide. Their
boundary endpoints therefore attach to lattice edges incident to both
the inner incident strand and the two external arms. These strand
attachments are on opposite sides of the incident
edge: the other possibility contradicts the winding number of
the circular separator completed along the common walk path.
This is the deterministic Jordan-curve conclusion in
\cite[final paragraphs of Lemma 3.17]{SSdiscrete}; its hypotheses
are the two coinciding crosscuts, a simple interface, and the
common walk avoiding it before zero, all obtained above.
Thus each external arm has merged with a distinct inner arm,
and there can be no additional annular interface crossing.

The total error is \(C\eps+o_{R/r}(1)+C(\eps)(1-c)^N\).
Choose the buffer error first and then \(R/r\) large.
Average over \(\Theta_r\) and the six possible incident edges
for the external-only statement.
\end{proof}

\subsection{Boundary clearance, survival, and nonreturn}

The survival and nonreturn arguments adapt
\cite[Lemma~3.26, especially equation (3.62)]{SSdiscrete}, using the
barrier and sampled-excursion estimates proved above.
For the boundary-attaching arguments in this subsection, fix
\(\lambda_*>0\) and assume the exterior Dobrushin data satisfy
\[
 h_\partial\in[\lambda_*,1/2]\text{ on }\partial_+,
 \qquad
 h_\partial\in[-1/2,-\lambda_*]\text{ on }\partial_-.
 \tag{GB}
\]
All constants may depend on \(\lambda_*\). This includes the reference
data \(\pm1/2\), and the candidate matched data \(\pm l\) with
\(l\geq c\) from Lemma~\ref{g:sign-trace}. No uniform assertion as
\(l\downarrow0\) is required. The source's boundary condition
\cite[condition \((\partial)\), Section 3.7]{SSdiscrete} is more
permissive, allowing small wrong-sign heights; we do not need or invoke
that additional Gaussian statement.

\begin{lemma}[Clearance from the original boundary]
\label{g:boundary-landing}
Let \(S\) be an independent full walk from \(v\), stopped only on
the original boundary at \(\tau_D\), and let \(\gamma\) be the
completed interface. Put \(d=\operatorname{dist}(v,\partial D)\) and
\[
 M=\{S_j:j<\tau_D,\ S_j\text{ is adjacent to }\gamma\}.
\]
Under \((\mathrm{GB})\), uniformly over domains, observers, and the
indicated exterior data,
\[
 \lim_{\eta\downarrow0}\limsup_{d\to\infty}
 \Prob\bigl(\operatorname{dist}(M,\partial D)<\eta d\bigr)=0.
 \tag{GB1}
\]
The distance of the empty set is infinity. In particular the assertion
holds for a terminal harmonic landing, or simultaneously for every
old-prefix landing of this walk that is not on \(\partial D\).
\end{lemma}

\begin{proof}
We adapt the full-walk argument of \cite[Lemma~3.23, especially
equations (3.53) and (3.55)]{SSdiscrete}. We first establish pure-arc
repulsion. For \(y\in\partial_+\), choose
any deterministic lower bound \(b\leq d_D(y,\partial_-)\).
At successive fixed-ratio scales between \(b\) and \(r\), take
the component of \(B(y,\rho)\cap D\)
accessible from \(y\), and its circular crosscut separating \(y\)
from \(\partial_-\). The clearance construction used in
Lemma~\ref{g:extra} supplies a barrier attached to the two positive
boundary pieces on either side of \(y\). Its endpoint pockets stay
on \(\partial_+\), since their diameter is smaller than \(b\).
Explore the interface only up to the next outer crosscut. Its exact
sign transcript leaves the inner barrier compatible, so
Proposition~\ref{g:barriers} gives a fixed conditional chance of
avoiding this crosscut and every smaller access to \(y\).
Iterating, then reflecting for \(y\in\partial_-\), gives constants
\(C,\zeta>0\) such that
\[
 \Prob\{d_D(y,\gamma)<r\}
 \leq C(r/b)^\zeta+o_d(1),                         \tag{GB2}
\]
when \(r,b\) are fixed positive multiples of \(d\) and \(r<b\).
The lattice error is uniform for these fixed multiples. This is the
deterministic crosscut argument of
\cite[Theorem 3.22]{SSdiscrete}, with its field input replaced by
our exterior version of Proposition~\ref{g:barriers}.

Fix \(r=\eta d\), and condition on \(\gamma\). At the first walk
visit to \(M\) at distance less than \(r\) from \(\partial D\),
the strong Markov property and the planar walk estimate give a
fixed probability of exiting \(D\) in the ball of radius \(2r+O(1)\)
about a closest boundary point. One can force the walk to surround
that point inside this ball, which forces the exit. Its continuation
and the adjacent interface edge then give an internal path of diameter
at most \(4r+O(1)\) from the exit \(Y=S_{\tau_D}\) to \(\gamma\).
Consequently
\[
 \Prob\{\operatorname{dist}(M,\partial D)<r\}
 \leq C\Prob\{d_D(Y,\gamma)<4r+O(1)\}.             \tag{GB3}
\]
Here \(Y\), unlike the interface landing, is independent of
\(\gamma\). Conditional on \(Y\), apply (GB2) whenever its internal
distance to the opposite arc exceeds \(s=\sqrt\eta\,d\), using
the deterministic lower bound \(b=s\) in that estimate. This gives
\(C\eta^{\zeta/2}+o_d(1)\).

It remains to bound exits within internal distance \(s\) of both
sign arcs. Delete all paths of diameter at most \(2s\) joining the
two arcs and retain the component of \(v\). Its new boundary lies
in two balls of radius \(4s\). To verify this elementary fact,
classify each deleted crosscut by which marked sign-change endpoint
it separates from \(v\). Any two crosscuts in the same class which
touch the retained component intersect: otherwise their separation
order prevents one of them from touching that component. Their
diameter bound puts the whole class in one ball. There are two
classes. An exit near both arcs is beyond this new boundary.
The walk hitting estimate for these two balls gives
\(C(s/d)^{1/2}=C\eta^{1/4}\). Combining with (GB3) proves (GB1).

Every nonexterior old-prefix landing is a vertex of \(M\), pathwise,
because the prefix is contained in \(\gamma\) and the full walk has
not yet exited \(D\). This proves the simultaneous assertion without
conditioning on any landing event or on any stopping time.
\end{proof}

\begin{lemma}[Clearance from the observer]\label{g:observer-landing}
Uniformly over the same bounded-data experiments,
\[
 \lim_{\eta\downarrow0}\limsup_{d\to\infty}
 \Prob\{\operatorname{dist}(v,\gamma)<\eta d\}=0.     \tag{GB4}
\]
This assertion does not require the lower bound in \((\mathrm{GB})\).
\end{lemma}

\begin{proof}
The nested shielding argument follows \cite[Lemma~3.24]{SSdiscrete},
using Proposition~\ref{g:barriers}. Expose the first interface approach
to the circles centered at \(v\)
of radii \(2^{-j-1}d\). Given its exact sign transcript at one such
circle, the circular clearance construction and
Proposition~\ref{g:barriers} give a fixed chance of protecting the
next smaller disk. All attachments are to the observed strands,
not to the original boundary. Thus failure through \(N\) scales
has probability at most \((1-p)^N\). Fix \(N\) first and let
\(d\to\infty\), so every used scale diverges; then let \(N\to\infty\).
This proves (GB4). It also excludes all terminal and old-prefix
landings within \(\eta d\) of \(v\), up to one lattice spacing.
\end{proof}

\begin{lemma}[Survival of a stopped harmonic landing]\label{g:survival}
Let \(T\) be any stopping time for the filtration generated by oriented
initial interface pieces, \(K=\gamma[0,T]\), and
\(d=\operatorname{dist}(v,\partial D)\). Let \(X_T\) be the first
landing of an independent full walk \(S\) on \(\partial D(K)\);
let \(Y\) be its landing for the completed interface. For every
\(\eps>0\) there are \(p>0\), \(d_0<\infty\) such that
\[
 \Prob\bigl(\Prob(Y=X_T\mid K,S)<p\bigr)<\eps
 \qquad(d>d_0),                                      \tag{G9}
\]
uniformly over the domains, observers, and stopping times under
\((\mathrm{GB})\).
\end{lemma}

\begin{proof}
Conditional on an exact stopped prefix \(K\), the field has the
ordinary strand-conditioned law: whether the rule stopped at this
prefix is already decided by that prefix. The independent walk
adds no field conditioning. Put \(U=D(K)\).
For small fixed \(\delta>0\), except on an event tending uniformly
to zero with \(\delta\), the following four conditions hold:
\[
\begin{split}
 &d_U(X_T,\hbox{opposite-sign side})\geq2\delta d,\\
 &\operatorname{diam} S[0,\tau_T]<\delta^{-1}d,\\
 &\operatorname{diam} S[t_*,\tau_T]<\delta d/4,
 \quad t_*=\inf\{t:d(S_t,\partial U)\leq\delta^2d\},\\
 &d(v,K)\geq\delta^{1/2}d.
\end{split}                                                   \tag{G10}
\]
The middle two statements are the planar walk hitting estimate
uniform in \(U\). The last statement follows from
Lemma~\ref{g:observer-landing} and \(K\subset\gamma\).
For the observer-separated stopping times needed in the driving
argument, the last condition can instead be imposed deterministically.

For the first condition delete all crosscuts of diameter at most
\(4\delta d\) joining the two sign sides. The new boundary of the
component containing \(v\) lies in two balls of radius \(8\delta d\).
Indeed each deleted crosscut separates \(v\) from one of the two
sign-change points; any two crosscuts for the same point which
touch the retained component must intersect. All such crosscuts
are consequently within twice their diameter of a fixed one.
Apply the walk hitting estimate to those two balls, using the last
condition. This also explains why arbitrary narrowing of \(U\)
does not invalidate (G10).

Fix records satisfying (G10), and suppose the landing is positive.
The deterministic clearance construction used in
Lemma~\ref{g:extra} gives a compatible crosscut \(\Upsilon_0\)
in \(B(X_T,\delta d)\setminus B(X_T,\delta d/2)\), joining the
positive boundary pieces on either side of \(X_T\) and separating
it from the negative side. If no hexagon centered at a vertex of the
stopped walk meets this crosscut,
avoidance of this crosscut protects the entire walk.
Otherwise add the boundaries of the walk hexagons on its
unprotected side and take the boundary facing the negative side.
Its diameter is at most \(2\delta^{-1}d\); its new part is at
distance at least \(\delta^2d/2\) from \(\partial U\), by the
terminal-walk condition. Indeed a contacting walk hexagon occurs
before \(t_*\), since the terminal walk piece lies in
\(B(X_T,\delta d/4)\), whereas the crosscut lies outside
\(B(X_T,\delta d/2)\). That contact is at distance at least \(\delta^2d-O(1)\)
from \(\partial U\); therefore the boundary-attached original
crosscut has diameter at least \(\delta^2d-O(1)\) as well. The old endpoints retain their pure
pockets. It is therefore a barrier with parameter \(c\delta^3\)
and scale \(2\delta^{-1}d\).

Proposition~\ref{g:barriers} gives probability at least \(p_\delta>0\)
that the future interface avoids the resulting barrier. On this
event it is adjacent to no vertex of the stopped walk, so \(Y=X_T\).
In the first branch \(\Upsilon_0\) can have bounded lattice diameter;
this is covered by the all-scale part of that proposition.
Reflect signs for a negative landing. Choose \(\delta\) to make
the failure probability of (G10) smaller than \(\eps\), and then
take \(d\) sufficiently large.
\end{proof}

\begin{proposition}[Nonreturn at stopped harmonic landings]
\label{g:nonreturn}
For every \(\eps>0\) there are \(s>0\) and \(d_0<\infty\) such that
\[
 \Prob\bigl(d(X_T,\gamma\setminus\gamma[0,T])<s d,\,
                  X_T\notin\partial D\bigr)<\eps
 \qquad(d>d_0),                                      \tag{G11}
\]
uniformly over all sign-filtration stopping times \(T\), with
\(d=\operatorname{dist}(v,\partial D)\), under \((\mathrm{GB})\).
\end{proposition}

\begin{proof}
Let \(\mathcal G=\sigma(K,S)\), \(A=\{Y=X_T\}\), and
\(q=\Prob(A\mid\mathcal G)\). Lemma~\ref{g:survival} implies,
for every \(\mathcal G\)-measurable event \(E\),
\[
 \Prob(E)\leq\eps_1+p^{-1}\Prob(E\cap A).              \tag{G12}
\]
Take \(E\) to be the event that the old prefix has more than two
crossings of an annulus about its walk landing. On \(A\), this
implies the corresponding event for the completed interface at
its completed harmonic landing. Lemma~\ref{g:extra}, with error
smaller than \(p\eps_1\), therefore removes all additional old
excursions from a small comparison window.

Exclude landings close to the old tip using the walk hitting bound.
Lemmas~\ref{g:boundary-landing} and \ref{g:observer-landing} exclude
landings close to the original boundary or observer, uniformly over
all prefixes. Choose the largest comparison window smaller than
one tenth of this retained clearance. This verifies the original
boundary and observer exclusions when Lemma~\ref{g:extra} is applied
after recentering at the landing.

In the remaining window exactly two incident old arms reach the
outer circle, and the old tip lies outside it. There are two
relevant complementary components. In each, select nested
compatible barriers attached to its pure old side. After each
failed shield, reveal the future arms up to that barrier; the
conditional field law is again an exact-strand law, so the next
shield succeeds with probability at least a fixed \(c>0\).
Failure through \(N\) scales has probability at most \((1-c)^N\).
When neither complementary component is traversed, the future
cannot enter the innermost Euclidean ball. The removal of
additional old excursions is what converts this internal-access
conclusion to Euclidean nonreturn.

Choose \(\eps_1\), obtain \(p\) from survival, choose the completed
excursion error below \(p\eps_1\), choose \(N\), and then let all
the finitely many retained scales tend to infinity in lattice
units. Finally choose their innermost relative radius \(s\).
The sum of the discarded probabilities can be made less than
\(\eps\), proving (G11).
\end{proof}

\begin{remark}[What the geometric result supplies]
The random mass of old-boundary accesses revisited by the future is
the conditional harmonic probability of the event in (G11).
Its expectation tends to zero, uniformly over stopping rules.
Consequently bounded harmonic boundary tests may discard this mass.
This is the estimate needed before a tower-property argument for a
stopped observable. Absolute continuity of access in each fixed domain
alone would not give its uniform form.
\end{remark}
